\documentclass[11pt]{article}

\usepackage[margin=1in]{geometry}
\usepackage{amsmath,amssymb,amsthm,mathtools,booktabs,array,tabularx}
\usepackage{microtype,xurl}
\usepackage{xcolor}
\usepackage[colorlinks=true,linkcolor=blue!55!black,
            citecolor=blue!55!black,urlcolor=blue!55!black]{hyperref}

\newtheorem{theorem}{Theorem}[section]
\newtheorem{proposition}[theorem]{Proposition}
\newtheorem{lemma}[theorem]{Lemma}
\newtheorem{corollary}[theorem]{Corollary}
\theoremstyle{definition}
\newtheorem{definition}[theorem]{Definition}
\theoremstyle{remark}
\newtheorem{remark}[theorem]{Remark}

\newcommand{\Q}{\mathbb Q}
\newcommand{\cC}{\mathcal C}
\newcommand{\cR}{\mathcal R}
\newcommand{\id}{\mathrm{id}}
\newcommand{\supp}{\operatorname{supp}}
\newcommand{\Cyc}{\operatorname{Cyc}}
\newcommand{\rank}{\operatorname{rank}}
\newcommand{\ellb}{\boldsymbol\ell}
\newcommand{\cert}[1]{\nolinkurl{#1}}

\title{The Gravity--Light Interaction Survives Cyclic Causal Reduction\\
through $\ell_3$}
\author{Asger Alstrup Palm\thanks{Corresponding author:
\texttt{asger@area9.dk}.}}
\date{Working draft, July 2026}

\begin{document}
\maketitle

\begin{abstract}
We transfer the classical gravity--Maxwell BV interaction on a compact
positive Berger-clock background from a complete $64$-row presentation to a
typed cyclic $36$-row retained complex.  The contraction carries the exact
odd pairing and the unary causal Green data.  In the repository's suspended
graded-symmetric convention the transferred binary operation has $1{,}474$
coefficients and the mixed ternary operation has $25{,}950$ coefficients.
The latter is not killed by reduction: it contains $7{,}614$ gravity-output
terms with two Maxwell inputs and $18{,}336$ Maxwell-output terms with one
Maxwell input.

The homological exchange channel is computed rather than omitted.  It
contains $342$ canonical full-complex coefficients in the only nonempty
channel, all supported on full row $38$, and the certified retained
projection annihilates that row.  A one-entry projection mutation exposes all
$342$ terms.  Thus the retained exchange contribution vanishes for a
structural reason while the direct $q_3$ contact survives.  All $36$
arity-three $L_\infty$ rows close.  Lowering $\ell_3$ with the retained typed
pairing independently reproduces all $25{,}662$ physical quartic
coefficients with zero cyclicity defects; changing the Maxwell pairing weight
from two to one produces $17{,}108$ exact defects on $14$ rows.
The remaining $288$ ghost/antifield coefficients also replay exactly under
the suspended-Darboux transpose, closing cyclicity of the full retained BV
tensor.

Together with the cyclic causal Cartan contraction for the
background-fixing helical generator $K_{\rm Berger}$, these calculations
establish
\[
 \boxed{\text{The gravity--light interaction survives cyclic causal
 reduction through }\ell_3.}
\]
The statement is classical and carries separate \textsc{local-algebraic}
and \textsc{lorentzian-causal} dependencies.  The causal claim concerns the
unary Green contractions and the two-sided-causal-hull Cartan primitives;
the Taylor operations themselves are support-local.  The retained complex is
not yet the Einstein-like/extra-Weyl residual model.  Its mixing table,
Hadamard products, QME restoration, and every quantum or particle
interpretation remain open.
\end{abstract}

\section*{AI-use disclosure and computational accountability}
OpenAI GPT-5 was used substantively for mathematical synthesis, verification
code, claim auditing, and manuscript drafting.  The human author directed the
programme and remains responsible for the claims, citations, and final text.
Every numerical count in the theorem is an exact sparse-coefficient count over
$\Q(\sqrt{10})$ and is tied to a deterministic certificate.  No floating
point calculation enters an identity, rank, cyclicity test, or mutation
witness.  Submission remains conditional on final human review.

\tableofcontents

\section{Question and answer}\label{sec:intro}

The free BV complex identifies which classes survive gauge reduction, but it
does not determine whether nonlinear gravity--matter vertices survive the
same reduction.  A transferred interaction can disappear in three distinct
ways: its direct contact term can project to zero; homological exchange terms
can cancel it; or cyclicity and the $L_\infty$ identities can fail after
compression.  Conversely, a nonzero tensor on a gauge-fixed carrier need not
descend to any physical branch until its residual projection is computed.

This paper answers the algebraic question through ternary order for the
gravity--Maxwell extension of the positive Berger-clock complex.  Write
\begin{equation}\label{eq:headline}
 \boxed{\text{The gravity--light interaction survives cyclic causal
 reduction through }\ell_3.}
\end{equation}
Here ``survives'' means that the explicitly transferred mixed ternary bracket
is nonzero and obeys the retained arity-three identity.  ``Cyclic'' means both
that the contraction is typed cyclic and that the complete physical quartic
form obtained by lowering $\ell_3$ is independently invariant under cyclic
transposition.  ``Causal reduction'' means that the unary contraction imports
the advanced/retarded Green homotopies and that the compatible
$K_{\rm Berger}$ Cartan primitives through arity three have support in the
two-sided causal hull.  It does \emph{not} mean that a time-ordered interacting
quantum theory has been constructed.

The theorem stops one rung before physical branch interpretation.  The
dynamical residual carrier must still be split into Einstein-like,
extra-Weyl, and Maxwell branches.  Separately, the even Weyl-square and odd
Pontryagin directions form a deformation/vertex basis; the odd topological
class is not a third particle branch.  This category separation is built into
the receiving contract for the next calculation.

\section{Complexes, conventions, and claim boundary}\label{sec:setup}

Let $(\cC_{64},q_1,\omega_{64})$ be the complete typed classical
gravity--clock--Maxwell BV complex at the frozen rational Berger fixture.
Its degree multiplicities are
\begin{equation}
 6,26,26,6\qquad\text{in degrees }-1,0,1,2.
\end{equation}
Let $(\cR_{36},\ell_1,\omega_{36})$ denote the retained typed carrier, with
degree multiplicities
\begin{equation}
 4,14,14,4.
\end{equation}
The exact contraction is
\begin{equation}\label{eq:sdr}
 \left(\cC_{64},q_1\right)
 \underset{\iota}{\overset{\pi}{\rightleftarrows}}
 \left(\cR_{36},\ell_1\right),
 \qquad
 \pi\iota=\id,
 \qquad
 q_1S+Sq_1=\id-\iota\pi,
\end{equation}
with side conditions and typed cyclic adjoint relations fixed by the carrier
certificate.  The coefficients lie in $\Q(\sqrt{10})$.  The Maxwell
field--equation pairing has absolute weight two, while the gravity weight is
one.  The signed odd-pairing ledger contains ten gravity entries of sign
$-1$ and four Maxwell entries of sign $+2$ on the physical
field--equation pairs.

The based classical BV vector field is known through cubic Taylor order,
\begin{equation}
 q=q_1+q_2+q_3+O(4),
\end{equation}
in the repository's suspended graded-symmetric factorial convention.  The
full tensors are action-derived, support-local, typed cyclic, and equivariant
under the background-fixing helical generator
\begin{equation}
 K_{\rm Berger}=D-\omega R.
\end{equation}
Raw cylinder translation $D$ is affine at the rotating background and is not
substituted for $K_{\rm Berger}$.

The result combines two dependency layers:
\begin{center}
\begin{tabularx}{0.96\textwidth}{>{\small}p{0.20\textwidth}XX}
\toprule
Layer & Certified content & Explicit limitation\\
\midrule
\textsc{local-algebraic}
 & exact $q_2,q_3,\ell_2,\ell_3$, identities and cyclicity
 & no residual branch interpretation\\
\textsc{lorentzian-causal}
 & unary Green homotopies and $K$-Cartan support through arity three
 & no Hadamard products or QME\\
\bottomrule
\end{tabularx}
\end{center}
Neither layer is silently promoted to \textsc{residual-transferred} or
\textsc{lorentzian-certified} quantum status.

\section{Ternary homological transfer}\label{sec:transfer}

The first two transferred Taylor operations are
\begin{align}
 \ell_1&=\pi q_1\iota,\label{eq:ell1}\\
 \ell_2&=\pi q_2(\iota,\iota).\label{eq:ell2}
\end{align}
Define the second Taylor component of the inclusion by
\begin{equation}\label{eq:I2}
 I_2=-S q_2(\iota,\iota).
\end{equation}
Then the ternary transfer formula is
\begin{equation}\label{eq:ell3}
 \ell_3
 =\pi q_3(\iota,\iota,\iota)
 +\sum_{\sigma\in\operatorname{Sh}(2,1)}
   \epsilon(\sigma)\,
   \pi q_2\bigl(I_2(\iota_{\sigma(1)},\iota_{\sigma(2)}),
                     \iota_{\sigma(3)}\bigr).
\end{equation}
The first term is the contact contribution.  The second line is the exchange
contribution $q_2Sq_2$.

\begin{lemma}[Exact exchange support]\label{lem:exchange}
For the mixed gravity--Maxwell ternary operation, the
gravity-outer/mixed-inner exchange channel contains $144$ outer/inner
coefficient pairs.  The three graded unshuffles produce $324$ signed
contributions and $342$ canonical full-complex PBW coefficients.  Every one
has full output row $38$.  The retained projection has no support on that row,
so this channel vanishes after projection.  The mixed-outer/gravity-inner and
mixed-outer/mixed-inner channels have no outer/inner pairs.
\end{lemma}

\begin{proof}
The exact consumer reconstructs the $96$ gravity and $12$ mixed coefficients
of $I_2$, including their full-row supports, before composing them with
$q_2$.  PBW normalization gives the stated counts.  Direct application of
$\pi$ annihilates the $342$ row-$38$ coefficients.  As a negative control,
adding the single mutant entry
\begin{equation}
 \pi_{\rm mut}(\text{retained row }0,\text{full row }38)=1
\end{equation}
exposes all $342$ coefficients.  Thus the zero is caused by the certified
projection and is not an empty table hard-coded by the consumer.
\end{proof}

\begin{corollary}[Contact form of the retained mixed bracket]
\label{cor:contact}
The retained mixed exchange contribution is zero and
\begin{equation}
 \ell_3^{\rm mixed}=\pi q_3^{\rm mixed}(\iota,\iota,\iota).
\end{equation}
This equality is a computed property of the declared contraction, not a
general truncation of the homotopy-transfer formula.
\end{corollary}

\section{The algebraic survival theorem}\label{sec:theorem}

\begin{theorem}[Gravity--light survival through $\ell_3$]
\label{thm:survival}
On the frozen positive Berger-clock gravity--Maxwell BV complex, the typed
cyclic contraction \eqref{eq:sdr} transfers the action-derived Taylor data
through ternary order with the following exact properties.
\begin{enumerate}
\item The retained mixed binary operation $\ell_2$ contains $1{,}474$
canonical coefficients.
\item The retained mixed ternary operation $\ell_3$ contains $25{,}950$
canonical coefficients on $18$ nonzero output rows.
\item All retained exchange sectors vanish for the support reason in
Lemma~\ref{lem:exchange}; the nonzero ternary operation is the transferred
contact term.
\item The retained arity-three identity closes on all $36$ rows:
\begin{equation}\label{eq:linf3}
 [\ell_1,\ell_3]+\frac12[\ell_2,\ell_2]=0.
\end{equation}
\item The mixed bracket has $7{,}614$ gravity-output coefficients with two
Maxwell inputs and $18{,}336$ Maxwell-output coefficients with one Maxwell
input.  Hence it is nonzero in both gravity--light directions.
\end{enumerate}
A localized one-coefficient mutation of $\ell_3$ creates two exact
arity-three defects.
\end{theorem}

\begin{proof}
Equations \eqref{eq:ell2}--\eqref{eq:ell3} are evaluated in exact sparse PBW
normal form over $\Q(\sqrt{10})$.  The independent consumer parses all
$59{,}598$ coefficients of the full mixed $q_3$ payload and reproduces all
$25{,}950$ retained contact coefficients with no missing, extra, or changed
terms.  Lemma~\ref{lem:exchange} evaluates every possible mixed exchange
channel.  Direct coefficient comparison of
\eqref{eq:linf3} gives zero defects on all rows.  The input-sector and
output-sector ledgers give the two nonzero counts in item 5.  Finally, the
mutation changes one retained coefficient and produces two normalized
defects, preventing acceptance by a consumer that merely assumes the expected
answer.
\end{proof}

The theorem is stronger than the statement that Maxwell fields occur in the
full action.  It says that their ternary coupling remains present after the
explicit gauge-fixed-to-retained homological reduction and after every
arity-three identity is imposed.  It is weaker than a residual particle
vertex: $\cR_{36}$ is retained, not yet minimal or decomposed into its
dynamical branches.

\section{Physical quartic cyclicity}\label{sec:cyclicity}

Lowering the output of $\ell_3$ with the odd pairing gives a four-linear
functional.  For physical degree-zero inputs, cyclic transposition of the
first input requires formal adjunction of every PBW derivative through the
invariant volume form.  If $w$ is a derivative word, the exact formal adjoint
distributes its letters over the other three slots with sign $(-1)^{|w|}$ and
the corresponding multinomial multiplicities.

\begin{proposition}[Independent physical quartic cyclicity]
\label{prop:cyclicity}
Every degree-zero physical coefficient of the lowered retained
$\ell_3$ satisfies exact cyclic transposition.  There are
\begin{equation}
 25{,}662=7{,}506_{\rm gravity\ output}
             +18{,}156_{\rm Maxwell\ output}
\end{equation}
such coefficients, and the defect count is zero on every row.  Replacing the
Maxwell field--equation weight two by weight one produces $17{,}108$ defects
on $14$ rows.
\end{proposition}

\begin{proof}
The verifier constructs both sides independently from the retained tensor,
the typed pairing, and exact PBW integration by parts.  The signed pairing
orientations are retained as a separate convention ledger; only their
absolute component multiplicities enter the physical field--equation
transpose.  Equality holds coefficientwise.  The weight mutation changes no
coefficient of $\ell_3$ itself and therefore tests the lowering convention
rather than the imported interaction.
\end{proof}

The remaining
\begin{equation}
 25{,}950-25{,}662=288
\end{equation}
coefficients are the ghost/antifield completion.  Their retained output rows
are
\[
 23,\ 24,\ 25,\ 26,\ 32,\ 33,\ 34.
\]
If $x$ is the transposed first input and $u$ is
the coordinate paired with the output, the suspended-Darboux sign is
\begin{equation}
 (-1)^{|x||u|+\epsilon_2(x)+\epsilon_2(u)},
 \qquad
 \epsilon_2(v)=\begin{cases}1,&\deg v=2,\\0,&\text{otherwise.}\end{cases}
\end{equation}
Here the degree-two rows are precisely the retained ghost-antifield
coordinates.

\begin{proposition}[Independent full-BV quartic cyclicity]
\label{prop:full-bv-cyclicity}
The complete $25{,}950$-coefficient lowered retained $\ell_3$ satisfies exact
cyclic transposition.  Within the $288$-coefficient completion, $120$
coefficients carry the positive transpose sign and $168$ carry the negative
sign.  The exact defect count is zero.  Omitting the degree-two polarization
term exposes $132$ defects on all seven completion output rows.
\end{proposition}

This proposition is a full-BV statement about the pinned classical retained
tensor.  It does not imply full BV local cohomology, a quantum master equation,
or a Lorentzian quantum theory.

\section{Causal compatibility through ternary order}\label{sec:causal}

The causal part of the headline is not inferred from finite PBW tables.  It
is supplied by the independently certified unary Green chain contractions
and the coupled cyclic Cartan theorem.  Let $\Lambda_{64,\pm}$ denote the
same-sided advanced and retarded unary homotopies.  Their chain identities,
causal supports, and formal-adjoint reversal are imported on the complete
carrier.  The local Taylor operations and the contraction intertwine
$K_{\rm Berger}$.

Let $\iota_K^{(1)}$ be the unary Cartan primitive.  The higher sources are
\begin{align}
 A_K^{(2)}&=[q_2,\iota_K^{(1)}]-L_K^{(2)},\label{eq:A2}\\
 A_K^{(3)}&=[q_3,\iota_K^{(1)}]+[q_2,\iota_K^{(2)}]-L_K^{(3)}.
 \label{eq:A3}
\end{align}
At the frozen background $L_K^{(2)}=L_K^{(3)}=0$.  The $L_\infty$ identities,
graded Jacobi identity, and $K$-equivariance give
\begin{equation}
 [q_1,A_K^{(2)}]=[q_1,A_K^{(3)}]=0.
\end{equation}
Raw causal primitives are cyclicized with the exact Reynolds projectors
\begin{equation}
 \Cyc_3=\frac{1+\tau+\tau^2}{3},
 \qquad
 \Cyc_4=\frac{1+\tau+\tau^2+\tau^3}{4}.
\end{equation}

\begin{theorem}[Cyclic causal Cartan compatibility]
\label{thm:cartan}
The coupled $64$-row complex and its typed $36$-row retained contraction obey
the $K_{\rm Berger}$ Cartan identities through arity three.  The higher cyclic
primitives satisfy
\begin{equation}\label{eq:support}
 \supp\bigl(\iota_K^{(n)}(f_1,\ldots,f_n)\bigr)
 \subseteq J\!\left(\bigcup_i\supp f_i\right),
 \qquad n\leq3.
\end{equation}
The Taylor operations are support-local and no spatial inverse or mode
projector enters the construction.
\end{theorem}

The higher cyclic primitives use the two-sided causal hull.  They are not
claimed to be separately retarded or advanced; the unary homotopies retain
that stronger one-sided property.  Theorem~\ref{thm:cartan} establishes
compatibility of the surviving interaction with the causal reduction.  It
does not construct interacting retarded products, a Hadamard state, or a
Lorentzian quantum master equation.

\section{The residual mixing table}\label{sec:mixing}

The next calculation determines what the surviving tensor means physically.
It requires a committed, content-addressed residual branch manifest rather
than an inferred basis.  The dynamical carrier must contain
\begin{equation}
 \cR_{\rm dyn}
 =\cR_{\rm Einstein}\oplus\cR_{\rm extra\ Weyl}
  \oplus\cR_{\rm Maxwell},
\end{equation}
together with exact inclusion, projection, pairing, parity, real structure,
and $K_{\rm Berger}$ weights.  The mixing entries are then
\begin{equation}\label{eq:mixing}
 M^{a}{}_{bcd}
 =\pi^{a}_{\rm dyn}\,
   \ell_3\bigl(\iota^{\rm dyn}_{b},
               \iota^{\rm dyn}_{c},
               \iota^{\rm dyn}_{d}\bigr).
\end{equation}
Every claimed zero must carry an exact support or cancellation witness, and
graded symmetry and cyclicity must be replayed in branch coordinates.

The deformation/vertex basis is separate:
\begin{equation}\label{eq:eo}
 e_{C^2}=\frac{[W_+^2]+[W_-^2]}{\sqrt2},
 \qquad
 o_{C\widetilde C}=\frac{[W_+^2]-[W_-^2]}{\sqrt2}.
\end{equation}
Its Euler--Lagrange map has the required rank pattern
\begin{equation}
 \rank E(e_{C^2})=1,
 \qquad
 E(o_{C\widetilde C})=0,
\end{equation}
with a separate Pontryagin transgression witness.  The odd topological
direction can be central or inert as a deformation class without becoming a
third dynamical particle branch.

At the time of this draft the strict projection consumer is ready, but the
authoritative branch-basis manifest has not been supplied.  Accordingly this
section freezes the question and the formula, not its answer.  The algebraic
survival theorem does not depend on that input.

\section{Interpretation and nonclaims}\label{sec:interpretation}

The result provides a precise first meaning of ``additional physics beyond
Einstein'' at nonlinear order: the retained classical interaction contains
both a gravitational response to two Maxwell inputs and a Maxwell response
to a gravity--Maxwell pair.  These channels survive exact cyclic transfer and
the ternary homotopy identity.  Once the branch manifest lands,
Equation~\eqref{eq:mixing} will decide whether the interaction closes on the
Einstein-like branch, couples it to the extra-Weyl branch, or acts only through
the deformation directions.

Several stronger conclusions do not follow.
\begin{enumerate}
\item The bracket changes no unary kinetic operator.  It therefore introduces
no new negative kinetic direction at this order, but this is not a unitarity
or positive-Hilbert-space theorem.
\item A nonzero retained bracket is not yet a photon or graviton scattering
amplitude.  Residual projection, causal state selection, and asymptotic or
relational observables are still required.
\item The odd Pontryagin direction is a topological deformation class, not a
one-particle state.
\item The vanishing one-particle cohomology of the separate free conformal
cylinder problem is not converted here into an interacting vacuum theorem.
\item No loop coefficient, anomaly cancellation, QME restoration, Hadamard
state, renormalized product, or quantum interaction has been constructed.
\item No arity-four or convergent all-orders $L_\infty$ or Cartan theorem is
claimed.
\end{enumerate}

The natural provisional interpretation is therefore a deformation/vertex
theory rather than a particle theory.  The mixing table is precisely the
next test of how much of that statement remains true after residual descent.

\section{Reproducibility and certificate map}\label{sec:repro}

The algebraic theorem is pinned to the classical hardening commit
\cert{e99d0c1d39490de5261fc6ca1dc2aeaa0d149655}.  Its principal artifacts are
\begin{center}
\begin{tabularx}{\textwidth}{>{\small}p{0.43\textwidth}X}
\toprule
Artifact & Role\\
\midrule
\cert{BERGER_COUPLED_K_CARTAN_THROUGH_ARITY_THREE}
 & full/retained cyclic causal $K$-Cartan theorem\\
\cert{BERGER_RETAINED_MIXED_ELL3_INDEPENDENT_ACCEPTANCE}
 & exact independent $64\to36$ transfer and mutation witnesses\\
\cert{BERGER_RETAINED_MIXED_ELL3_PHYSICAL_CYCLICITY}
 & independent physical quartic cyclicity and pairing-weight mutation\\
\cert{BERGER_RETAINED_MIXED_ELL3_FULL_BV_CYCLICITY}
 & independent ghost/antifield completion and full retained BV cyclicity\\
\cert{BERGER_RESIDUAL_MIXED_ELL3_BRANCH_PROJECTION_READINESS}
 & fail-closed input contract for the mixing table\\
\bottomrule
\end{tabularx}
\end{center}

The scoped reproduction commands are
\begin{verbatim}
PYTHONPATH=. python3 \
  d_quotient_classical/backreacted_clock/\
  berger_coupled_k_cartan_through_arity_three.py --check --guards
PYTHONPATH=. python3 \
  d_quotient_classical/backreacted_clock/\
  verify_berger_coupled_k_cartan_through_arity_three.py
PYTHONPATH=quantum-weyl python3 -m \
  transfer.berger_retained_mixed_ell3_acceptance_certificate --check
PYTHONPATH=quantum-weyl python3 -m \
  transfer.verify_berger_retained_mixed_ell3_acceptance
PYTHONPATH=quantum-weyl python3 -m \
  transfer.berger_retained_mixed_ell3_physical_cyclicity_certificate --check
PYTHONPATH=quantum-weyl python3 -m \
  transfer.verify_berger_retained_mixed_ell3_physical_cyclicity
PYTHONPATH=quantum-weyl python3 -m \
  transfer.berger_retained_mixed_ell3_ghost_antifield_cyclicity_certificate --check
PYTHONPATH=quantum-weyl python3 -m \
  transfer.verify_berger_retained_mixed_ell3_ghost_antifield_cyclicity
PYTHONPATH=quantum-weyl python3 -m \
  transfer.verify_berger_residual_ell3_branch_projection_readiness
\end{verbatim}

The companion machine-readable claim map records the dependency hashes,
positive claims, explicit nonclaims, and the fact that this manuscript is a
working draft rather than a theorem freeze.

\section{Outlook}\label{sec:outlook}

Paper~11 can now develop on two rails.  The first is editorial: sharpen the
abstract transfer theorem, extract compact coefficient tables to a
supplement, and obtain specialist review of the cyclic-sign and causal-support
conventions.  The second is mathematical: import the residual branch basis
and compute Equation~\eqref{eq:mixing}.  A nonzero
Einstein/extra-Weyl entry would turn the additional Weyl physics into an
explicit interaction channel; a block-diagonal table would establish
nonlinear closure of the Einstein-like direction through this order; and a
central topological action would isolate the vertex theory from the
dynamical carrier.

The algebraic theorem does not need to wait for that verdict.  What the
mixing table changes is its physical interpretation, not the existence of the
surviving cyclic gravity--light bracket.

\begin{thebibliography}{99}

\bibitem{BV}
I.~A. Batalin and G.~A. Vilkovisky,
\emph{Gauge algebra and quantization},
Phys. Lett. B \textbf{102} (1981), 27--31.

\bibitem{HT}
M.~Henneaux and C.~Teitelboim,
\emph{Quantization of Gauge Systems},
Princeton University Press, 1992.

\bibitem{Kajiura}
H.~Kajiura,
\emph{Noncommutative homotopy algebras associated with open strings},
Rev. Math. Phys. \textbf{19} (2007), 1--99.

\bibitem{HomotopyTransfer}
A.~S. Arvanitakis, O.~Hohm, C.~Hull, and V.~Lekeu,
\emph{Homotopy transfer and effective field theory I: tree level},
arXiv:2007.07942.

\bibitem{FiorenzaManetti}
D.~Fiorenza and M.~Manetti,
\emph{$L_\infty$ algebras, Cartan homotopies and period maps},
arXiv:math/0605297.

\bibitem{LInfinityIntro}
A.~Kraft and J.~Schnitzer,
\emph{An introduction to $L_\infty$-algebras and their homotopy theory},
arXiv:2207.01861.

\bibitem{GreenComplexes}
M.~Benini, G.~Musante, and A.~Schenkel,
\emph{Green hyperbolic complexes on Lorentzian manifolds},
Commun. Math. Phys. (2023),
\href{https://doi.org/10.1007/s00220-023-04807-5}
{doi:10.1007/s00220-023-04807-5}.

\end{thebibliography}

\end{document}
